% year: תשפ"ד
% type: מבחן
% semester: א
% moed: מיוחד
% number: 3
% points: 14
% categories: func-analysis, derivatives
% source: exams/מבחנים/תשפד/מועד ג תשפד.pdf

\begin{question}
חשבו תחומי מינימום ומקסימום מוחלטים של הפונקציה
\[ f(x)=(x+2)\cdot e^x \]
בתחום הגדרתה, או הסבירו מדוע אינם קיימים.
\end{question}

\begin{hint}
מצאו את תחומי העלייה והירידה בעזרת $f'$, ובדקו את התנהגות הפונקציה כאשר $x\to\pm\infty$.
\end{hint}

\begin{solution}
\textbf{תחום הגדרה:} כל $\R$; $f$ גזירה בכל נקודה.

\textbf{נגזרת:} לפי כלל המכפלה
\[ f'(x)=e^x+(x+2)e^x=(x+3)e^x. \]
כיוון ש-$e^x>0$, סימן $f'$ הוא סימן $x+3$: $f'<0$ ב-$(-\infty,-3)$ ו-$f'>0$ ב-$(-3,\infty)$. לכן (לפי מסקנה ממשפט לגרנז') $f$ יורדת ממש ב-$(-\infty,-3]$ ועולה ממש ב-$[-3,\infty)$.

\textbf{מינימום מוחלט:} מהמונוטוניות, לכל $x\le-3$ מתקיים $f(x)\ge f(-3)$, ולכל $x\ge-3$ מתקיים $f(x)\ge f(-3)$. לכן $x=-3$ היא נקודת מינימום מוחלט, וערך המינימום הוא
\[ f(-3)=(-3+2)e^{-3}=-\frac1{e^3}\approx-0.0498. \]

\textbf{מקסימום מוחלט:} $\lim_{x\to\infty}(x+2)e^x=+\infty$, ולכן הפונקציה אינה חסומה מלעיל ואין לה מקסימום מוחלט. (לשלמות: $\lim_{x\to-\infty}(x+2)e^x=\lim_{x\to-\infty}\frac{x+2}{e^{-x}}=0$ לפי לופיטל, כך שהפונקציה שואפת ל-$0^-$ משמאל, אך מימין אינה חסומה.)

\textbf{תשובה:} מינימום מוחלט $-e^{-3}$ המתקבל ב-$x=-3$; מקסימום מוחלט אינו קיים כי $f(x)\to+\infty$ כאשר $x\to\infty$.
\end{solution}
